\setcounter{secnumdepth}{-1}
\chapter{Згадані задачі}
\label{append:1}
\setcounter{equation}{0}
\renewcommand{\theequation}{\arabic{equation}}

\noindent Задача 3.1. Довести, що булева функція є алгебраїчно виродженою тоді й тiльки тоді, коли вона має ненульовий 
несуттєвий вектор.

Задача 3.2. Довести, що булева функція вiд $n$ змінних є $s$-вимірною тоді й тiльки тоді, коли вона має не менше 
нiж $n-s$ лінійно незалежних несуттєвих векторів, $s \in \overline{0, n-1}$.

Задача 3.3. Нехай $f$ є $s$-вимірною, але не $(s-1)$-вимірною функцією вiд $n$ змінних, $s \in \overline{1, n-1}$. 
Довести, що існує булева $n \times s$ - матриця $A$ рангу $s$ та булева функція $g$ вiд $s$ змінних, яка не має 
ненульових несуттєвих векторів, такі, що $f(x) = g(xA)$ для кожного $x \in V_{n}$. Чи визначається така пара $(A, g)$ 
для функції $f$ однозначно?

Задача 3.21. Довести, що булева функція вiд $n$ змінних є $s$-вимірною, $s \in \overline{0, n-1}$, тоді й тiльки 
тоді, коли множина $\left\{\alpha \in V_{n} \; : \; \hat{f}(\alpha) \neq 0\right\}$ породжує у векторному просторі 
$V_{n}$ підпростір розмірності не вище за $s$.

Задача 3.24. Нехай $f$ -- булева функція вiд $n$ змінних,
\begin{equation*}
    L_{f} = \left\{\alpha \in V_{n} \vert \forall x \in V_{n} \; : \; f(x \oplus a) = f(x)\right\}
\end{equation*}
Визначимо булеву функцію $h$ вiд $n$ змінних, вважаючи, що 
$h(x) = 1 \Leftrightarrow \hat{f}(x) \neq 0, \, x \in V_{n}$. Довести, що $\alpha \in V_{n}\setminus\left\{0\right\}$ 
належить $L_{f}$ тоді й тiльки тоді, коли $C_{h}(\alpha) = C_{h}(0)$. Переконайтесь, що $L_{f} = \left\{0\right\}$, 
якщо $\Vert h \Vert > 2^{n-1}$
